\zhbox{54}{64}{235}{670}{10}{\textbf{强制连续导通模式（FCCM）}\par
MP2984 可工作于 FCCM。降压导通时间（$t_{ON\_BUCK}$）和升压关断时间（$t_{OFF\_BOOST}$）由内部电路决定，使器件根据 $V_{IN}/V_{OUT}$ 比值以固定开关频率（$f_{SW}$）工作。负载减小时，输入电流（$I_{LIMIT}$）下降，可能使 $I_L$ 在 SWD 导通期间从 $V_{OUT}$ 流向 $V_{IN}$ 而变为负值。FCCM 强制 $I_L$ 以固定 $f_{SW}$ 工作，从而获得较低的输出纹波。\par
\textbf{开关电流限制}\par
MP2984 在控制环路中通过电流检测电阻（$R_{SENSE1}$）检测低侧 MOSFET 电流，在每个开关周期内于降压模式提供谷值限流，于升压模式提供峰值限流。在降压模式下，只有 $I_L$ 降至谷值限流值，下一周期才开始，因此触发谷值限流时器件可能降低频率。降压限流阈值固定为133 mV，升压限流阈值固定为150 mV。例如 $R_{SENSE1}=5$ m$\Omega$ 时，降压限流为133 mV/5 m$\Omega=26.6$ A，升压限流为150 mV/5 m$\Omega=30$ A。\par
开关电流限值可通过外部检测电阻配置。SWB 与 SWC 的电流信号在内部有约180 ns消隐时间，以增强抗噪能力。\par
发生过流时，MP2984 采用逐周期限流工作。器件可根据 OCP\_MODE 启动打嗝保护或锁断保护。\par
在打嗝模式下，软启动结束后，若 $V_{FB}<60\%V_{REF}$ 且触发开关限流，器件关断。在由 SS 电容放电周期配置的固定关断时间后恢复。\par
在锁断模式下，软启动结束后，若 $V_{FB}<60\%V_{REF}$ 且触发开关限流，器件关断，重新上电后恢复。\par
若禁用打嗝及锁断保护，器件持续以逐周期限流开关工作。软启动期间，打嗝及锁断保护被屏蔽。}
\zhbox{324}{64}{235}{620}{10}{\textbf{平均电流限制}\par
IAVGP 与 IAVGN 引脚检测输出电流。可在输出线上连接电流检测电阻（$R_{SENSE2}$），以控制平均输出电流限值。当检测信号超过限流参考电压时，内部 EA 将 $V_{SS}$ 拉低。最终，$V_{SS}$ 替代 $V_{REF}$ 控制 COMP，COMP 限制电感电流，使向输出传递的能量减少。SS 引脚使 $V_{OUT}$ 保持低电平，直到平均负载电流下降。\par
若开关电流由平均电流限制调节，且没有触发逐周期限流，则即便达到平均电流限值，MP2984 也不执行打嗝或锁断保护。这使器件能够保持恒流充电。\par
ILIM 位提供8种可选电流限值，外部 IPWM 引脚提供高分辨率限值调节。\par
若 IPWM 为高或悬空，平均输出电流限值由 ILIM 位设置与 $R_{SENSE2}$ 共同决定。\par
另一种设置方法是将 ILIM 位设定值与 IPWM 占空比相乘。ILIM 位为111且由 IPWM 引脚配置限流时，阈值 OCP\_LIMIT 可由式(3)计算：\par
\resizebox{\linewidth}{!}{$\mathrm{OCP\_LIMIT}\,(\mathrm{mV})=65.8\,\mathrm{mV}\times\mathrm{DUTY}+4.49\,\mathrm{mV}\quad(3)$}\par
其中 OCP\_LIMIT 是由 IPWM 配置的平均负载电流限值（ILIM 位为111），DUTY 为 IPWM 输入信号占空比，范围0--0.9。\par
实际输出电流限值（A）由 $\mathrm{OCP\_LIMIT}/R_{SENSE2}$ 决定。\par
例如，IPWM 占空比为48\%时，最终平均输出限流阈值约36 mV。IPWM 信号频率范围为5--100 kHz，建议采用20 kHz。对于 USB PD 应用，建议 $R_{SENSE2}=10$ m$\Omega$。}
\zhbox{324}{691}{234}{29}{10}{\textbf{过载保护（OLP）和短路保护（SCP）}}
